\scanpage{19}
mieux, si la \uline{famille} des $\dot T_i\subset\dot T$ au dessus
de $Y$ est \deleted{globalement} loc. triviale. On étend cette
définition au cas de $X_i\subset X$, et d'un
voisinage \deleted{tubulaire}\inserted{conique} $\bar T$ de $Y$ dans $X$ :
on dit
il est équisingulier p.r. à $(X_i)$, si $\bar T$
l'est par rapport à $(X_i\cap\bar T\subset\bar T)_i$.
Cela
implique \inserted{\unsure{f ?}} que si $\forall i$, $X_i\cap Y$ est une
partie à la fois ouverte et fermée de $Y$
\unsure{alors}, si $Y$ est connexe, que \deleted{$X_i\cap Y$}
l'on ait l'alternative
\[
(14)\qquad X_i\cap\bar T=\varnothing\quad\underline{\text{ou}}\quad X_i\supset Y
\]
\margin{on dit qu'il est \uline{strictement équisingulier}
si de plus il est \uline{strictement} \unsure{compatible avec tous}
les $X_i$.
Si $\bar T$ équisingulier p.r. à $(X_i)$,
$\lambda\bar T$ l'est strictement pour $\lambda\in]0,1[$.}

\deleted{\textcircled{\unsure{5}}}
On dit aussi que $X,(X_i)_{i\in I}$ est
\inserted{(globalement)} équisingulier
le long de $Y$, s'il existe un voisinage conique
\deleted{\unsure{de}} $\bar T$ de $Y$ dans $X$, qui est équisingulier
p.r. \deleted{à} $(X_i)$ --- auquel cas les $\lambda\bar T$
$(\lambda\in]0,1[)$ sont $=$ strictement équisinguliers.
\margin{18}
